% Created 2026-06-27 Sat 13:15
% Intended LaTeX compiler: pdflatex
\documentclass[11pt,letterpaper]{article}
\usepackage[utf8]{inputenc}
\usepackage[T1]{fontenc}
\usepackage{graphicx}
\usepackage{longtable}
\usepackage{wrapfig}
\usepackage{rotating}
\usepackage[normalem]{ulem}
\usepackage{amsmath}
\usepackage{amssymb}
\usepackage{capt-of}
\usepackage{hyperref}
\usepackage{amsmath}
\usepackage{amssymb}
\usepackage{geometry}
\geometry{left=1.2in, right=1.2in, top=1.0in, bottom=1.0in}
\usepackage{fancyhdr}
\usepackage{booktabs}
\usepackage{array}
\usepackage{parskip}
\usepackage{microtype}
\usepackage{tcolorbox}
\tcbuselibrary{skins,breakable}
\newenvironment{plainbox}{\begin{tcolorbox}[colback=blue!3!white,colframe=blue!40!white,title={\small\textit{In plain language}},breakable]}{\end{tcolorbox}}
\newenvironment{openbox}{\begin{tcolorbox}[colback=orange!5!white,colframe=orange!60!black,title={\small\textit{Open calculation}},breakable]}{\end{tcolorbox}}
\newenvironment{defbox}{\begin{tcolorbox}[colback=green!3!white,colframe=green!45!black,title={\small\textit{Foundational definition}},breakable]}{\end{tcolorbox}}
\PassOptionsToPackage{hidelinks}{hyperref}
\setlength{\parindent}{0pt}
\setlength{\parskip}{6pt}
\fancypagestyle{plain}{\fancyhf{}\fancyfoot[C]{\thepage}\renewcommand{\headrulewidth}{0pt}}
\fancypagestyle{body}{%
\fancyhf{}%
\fancyhead[C]{\small Document 3 of 5 $\cdot$ DOI: 10.5281/zenodo.18927154 $\cdot$ CC BY 4.0}%
\fancyfoot[C]{\thepage}%
\renewcommand{\headrulewidth}{0.4pt}}
\author{Bradley K. DePuy, Independent Researcher}
\date{2026}
\title{The Lagrangian of the Negotiation Field}
\hypersetup{
 pdfauthor={Bradley K. DePuy, Independent Researcher},
 pdftitle={The Lagrangian of the Negotiation Field},
 pdfkeywords={},
 pdfsubject={},
 pdfcreator={Emacs 29.3 (Org mode 9.6.15)}, 
 pdflang={English}}
\begin{document}

\pagestyle{plain}
\begin{titlepage}
\vspace*{3cm}
\begin{center}
  {\large\bfseries THE MEDIUM --- AN ATTEMPTED DESCRIPTION}\\[1em]
  {\LARGE\bfseries The Lagrangian of the Negotiation Field:}\\[0.4em]
  {\large Connecting a Pre-Geometric Substrate to Quantum Field Theory}\\[0.3em]
  {\normalsize Version 2 --- with Foundational Iota Scale Definition}\\[2em]
  {\normalsize Bradley K. DePuy, Independent Researcher}\\[0.4em]
  {\normalsize ORCID: \texttt{0009-0001-0328-2299}}\\[0.4em]
  {\normalsize Zenodo DOI: 10.5281/zenodo.18927154}\\[0.4em]
  {\normalsize github.com/bkdepuy/The-Medium}\\[0.4em]
  {\normalsize 2026}
\end{center}
\end{titlepage}
\pagestyle{body}

\section{Abstract}
\label{sec:org9458480}

This paper derives the Lagrangian of the negotiation field within The Medium and establishes its connections to quantum field theory, the Standard Model, and quantum gravity. The foundational definition --- iota size \(1 =\) Planck length --- is adopted from Document 2 (Unified Derivation v3). In Planck units, the Lagrangian takes a clean form: the kinetic term is the dimensionless cycling rate of the negotiation field, and the potential term is the graph Laplacian quadratic form. The Planck constant \(\hbar\) is identified as the projection seam between the dimensionless substrate action and the observer's dimensional units. Note: the Planck-unit route \(\hbar = E_\text{Planck} \cdot t_\text{Planck}/2\pi\) is algebraically circular --- \(E_\text{Planck} \cdot t_\text{Planck} = \hbar\) by definition of Planck units. The non-circular approach (deriving \(\hbar\) from proton mass, proton radius, and \(c\)) is discussed in Document 2. The symmetry group of the negotiation Lagrangian under complex negotiation states is at minimum U(1) × U(2); whether additional structure produces SU(3) specifically is the open derivation. The path integral over negotiation field configurations connects directly to quantum field theory. Four open calculations are identified: the correct symmetry group of the triangle Lagrangian, the Barbero-Immirzi correction to the eigenvalue spectrum, the Feynman rules from the topology change term, and the precise derivation of \(\hbar\).

\begin{defbox}
\textbf{Foundational definition (adopted from Document 2):}\\[4pt]
Iota size $1 =$ Planck length $= 1.616 \times 10^{-35}$ m\\[4pt]
In Planck units: $E_\text{Planck} = 1$, $t_\text{Planck} = 1$, $l_\text{Planck} = 1$, $\hbar = 1/(2\pi)$.\\[4pt]
The Lagrangian is dimensionless at the substrate level --- a pure ratio of negotiation energies in iota-scale units.
\end{defbox}

\section{The Lagrangian in Planck Units}
\label{sec:orgee2c750}

\begin{plainbox}
The Lagrangian is the master equation of this description --- one equation from which all physics flows. In Planck units it has a remarkably simple form: the square of how fast the negotiation is changing, minus the square of how much unresolved tension exists. The first term drives mass. The second term selects which patterns are stable.
\end{plainbox}

\subsection{The Field Variable}
\label{sec:org474a340}

The negotiation field \(N_{ij}\) is defined on every edge \((i,j)\) of the iota graph. In Planck units it is dimensionless --- the boundary resolution expressed as a ratio in iota-scale units. In the continuum limit it becomes \(N(x,t)\), a scalar field on emergent spacetime.

\subsection{The Lagrangian Density}
\label{sec:org9947173}

In Planck units (\(l_\text{Planck} = t_\text{Planck} = E_\text{Planck} = 1\)):

\begin{equation}
\boxed{\mathcal{L} = \frac{1}{2}\!\left(\frac{\partial N}{\partial t}\right)^{\!2} - \frac{1}{2}\,N^\top L\, N}
\end{equation}

The kinetic term \(\tfrac{1}{2}(\partial N/\partial t)^2\) is the energy of negotiation cycling in iota units. The potential term \(\tfrac{1}{2}(N^\top L N)\) is the closure penalty in iota units. Both are dimensionless ratios. The Lagrangian itself is dimensionless at the substrate level.

Restoring physical units:

\begin{equation}
\mathcal{L}_\text{SI} = \frac{E_\text{Planck}^2}{l_\text{Planck}^3}\left[\frac{1}{2}\!\left(\frac{\partial N}{\partial t}\right)^{\!2} - \frac{1}{2}\,N^\top L\, N\right]
\end{equation}

The prefactor \(E_\text{Planck}^2/l_\text{Planck}^3\) is the Planck energy density --- the natural unit of Lagrangian density.

\subsection{The Action}
\label{sec:org03c5ac4}

\begin{equation}
S = \int \mathcal{L}\, dt, \qquad [S] = E_\text{Planck} \cdot t_\text{Planck} = \hbar \cdot 2\pi
\end{equation}

The action has units of \(\hbar \cdot 2\pi\) --- which in Planck units is simply \(2\pi\). The phase factor in the path integral \(e^{iS/\hbar}\) becomes \(e^{i \cdot 2\pi \cdot S_\text{Planck}}\), a pure phase. Quantum mechanics is the interference of pure phases generated by dimensionless action ratios.

\section{Equations of Motion and the Wave Equation}
\label{sec:orga3012e2}

\begin{plainbox}
The equation that comes out of the Lagrangian is the wave equation --- confirming that negotiation modes are waves. In Planck units, the wave speed is $1 = c$. All massless negotiation modes propagate at the speed of light. Massive particles are standing waves whose cycling gives them inertia.
\end{plainbox}

Applying the Euler-Lagrange equation to \(\mathcal{L}\):

\begin{equation}
\frac{\partial^2 N}{\partial t^2} = -L \cdot N
\end{equation}

In the continuum limit: \(\partial^2 N/\partial t^2 = \nabla^2 N\) (wave speed \(= 1 = c\) in Planck units).

Standing wave solutions on a closed graph of \(n\) nodes:

\begin{equation}
N_k(t) = A_k \cdot v_k \cdot \cos(\omega_k t + \phi_k), \qquad \omega_k = \sqrt{\lambda_k}
\end{equation}

For the triangle: \(\omega = \sqrt{3}\) in Planck units. The energy \(E = \hbar\omega = \hbar\sqrt{3}\). Converting to MeV and applying the Barbero-Immirzi correction that shifts \(\omega\) from \(\sqrt{\lambda}\) to \(\lambda\) gives \(E = 3\varepsilon = 938.3\) MeV.

\begin{openbox}
\textbf{Open calculation 1 --- harmonic to linear spectrum}\\[4pt]
The Lagrangian gives $\omega_k = \sqrt{\lambda_k}$ (harmonic spectrum). The observed proton mass corresponds to $\omega_k = \lambda_k$ (linear spectrum). The Barbero-Immirzi parameter $\gamma \approx 0.2375$ sets the minimum area quantum and is expected to modify the dispersion relation from harmonic to linear at high energy density. Deriving this modification precisely from $\gamma \cdot l_\text{Planck}^2$ is the primary open calculation of the series.
\end{openbox}

\section{Symmetries in Planck Units}
\label{sec:orgd25a7d5}

\begin{plainbox}
The discrete automorphism group of the triangle graph \(K_3\) is \(S_3\) --- the symmetric group on three elements (six permutations). This is what the topology alone gives. When the negotiation states are taken to be complex-valued, the Lagrangian acquires additional continuous symmetries. For three complex edge states \(N \in \mathbb{C}^3\), the group of unitary transformations preserving the quadratic form \(N^\dagger L N\) is a subgroup of U(3). Whether this subgroup is precisely SU(3) depends on the specific constraint structure of the closure condition, and whether the physical identification of the triangle edges with color degrees of freedom is correct. The claim that ``SU(3) falls out of the triangle'' is directionally plausible but not yet a complete derivation --- it requires establishing that the complex negotiation states transform under SU(3) specifically, and that this corresponds to the QCD color symmetry.
\end{plainbox}

\subsection{U(1) Symmetry --- Electromagnetism}
\label{sec:orgf577ea2}

Rotating all negotiation phases simultaneously by angle \(\theta\) in Planck units:
\(N \to e^{i\theta} N\) leaves \(\mathcal{L}\) unchanged. This U(1) symmetry, when made local (\(\theta \to \theta(x,t)\)), requires a gauge field. In Planck units, the gauge field is massless, propagating at \(c\). This is the photon.

\subsection{SU(3) Symmetry --- Strong Force}
\label{sec:org5bcdb11}

The automorphism group of the triangle \(K_3\) is \(S_3\). When negotiation states are complex-valued in Planck units, the Laplacian \(L\) has eigenvalues \(\{0, 3, 3\}\). The unitary transformations preserving \(N^\dagger L N\) must preserve the \(\lambda=0\) eigenspace and the \(\lambda=3\) eigenspace separately; the correct symmetry group is U(1) × U(2), not SU(3). SU(3) would require \(L \propto I\), which holds only for the trivial graph.

Whether additional structure --- a cubic invariant from the antisymmetric tensor, or a different action of the symmetry on the edge states --- could produce SU(3) specifically is the open derivation. If it can be established, its eight generators would correspond to eight gluon states and the three color charges would be the three complex components of \(N\). That identification is a structural candidate, not yet a result.

\section{The Path Integral in Planck Units}
\label{sec:orgf42dd73}

\begin{plainbox}
In Planck units, the path integral sums over all possible negotiation histories, each weighted by a pure phase $e^{i \cdot 2\pi \cdot S}$ where $S$ is the action in units of $\hbar$. The most probable histories are those where the action is stationary. Stationary action means the negotiation cycles efficiently while staying close to closure. The proton is the most probable three-iota negotiation history.
\end{plainbox}

\subsection{The Partition Function}
\label{sec:org87946fb}

\begin{equation}
Z = \int \mathcal{D}N \cdot e^{iS/\hbar}, \qquad \hbar = \frac{1}{2\pi}\ \text{(Planck units)}
\end{equation}

\begin{equation}
Z = \int \mathcal{D}N \cdot e^{i \cdot 2\pi S}
\end{equation}

In Planck units, \(S\) is dimensionless and \(e^{i2\pi S}\) is a pure phase on the unit circle. The path integral is the sum of unit-circle phases over all negotiation histories. Constructive interference selects the stable particles.

\subsection{\(\hbar\) as a Derived Quantity}
\label{sec:org36472a2}

In Planck units, \(\hbar = 1/(2\pi)\) is not an independent constant. It is the conversion factor between the dimensionless Planck-unit action and the physical action in J\(\cdot\)s:

\begin{equation}
\hbar = \frac{E_\text{Planck} \cdot t_\text{Planck}}{2\pi}
\end{equation}

The standard Planck units are defined via \(E_\text{Planck} = \sqrt{\hbar c^5/G}\) and \(t_\text{Planck} = \sqrt{\hbar G/c^5}\), so \(E_\text{Planck} \cdot t_\text{Planck} = \hbar\) by construction. If the iota scale is defined geometrically as the minimum negotiation extent independent of \(\hbar\), then \(E_\text{Planck}\) and \(t_\text{Planck}\) must be derived from \(l_\iota\), \(G\), and \(c\) without invoking \(\hbar\). That derivation requires the substrate strain field equations, which are open work. Until then, \(\hbar\) in Planck units (\(\hbar = 1/2\pi\)) is a unit convention, not a result.

\section{Graph Topology Change and Particle Creation}
\label{sec:org9ee436c}

\begin{plainbox}
Particles are created when enough energy concentrates locally to open a new set of negotiation edges --- forming a new closed graph. In Planck units, the energy threshold for creating a proton is $3\varepsilon \approx 938$ MeV, which is ${\sim}10^{-20}$ in iota-scale units. This is why particle creation is rare at ordinary energies but routine in high-energy collisions.
\end{plainbox}

\subsection{The Extended Lagrangian}
\label{sec:org183855a}

\begin{equation}
\mathcal{L}_\text{full} = \mathcal{L}_\text{field} + \mathcal{L}_\text{topology}
\end{equation}

The topology change term in Planck units:

\begin{equation}
\mathcal{L}_\text{topology} = -\varepsilon_\text{edge} \sum_{ij} |\dot{h}_{ij}|
\end{equation}

where \(\varepsilon_\text{edge}\) is the energy cost of creating or dissolving one negotiation edge in Planck units. For a proton (3 nodes, 3 edges), the total creation energy is \(3\varepsilon_\text{node} + 3\varepsilon_\text{edge} = 938.3\) MeV.

\begin{openbox}
\textbf{Open calculation 2 --- Feynman rules from topology change}\\[4pt]
The topology change Lagrangian generates particle creation and annihilation. Deriving the Feynman rules --- the precise amplitudes for each process in Planck units --- would establish The Medium as a fully calculational description. The key quantity is $\varepsilon_\text{edge}$ in iota units, derivable from the geometric mean formula modified by the topology of the specific edge being created.
\end{openbox}

\section{Gravity and the Cosmological Constant}
\label{sec:orgc113f6e}

\begin{plainbox}
In Planck units, gravity is the curvature of the negotiation graph --- how much the local density of iota negotiations deviates from the average. The cosmological constant is the minimum negotiation energy density of the universe in Planck units --- one Hubble-floor energy per Hubble volume. It would follow from the iota grain-to-horizon ratio if the Hamiltonian floor condition can be derived --- that derivation does not yet exist.
\end{plainbox}

\subsection{Gravitational Coupling in Planck Units}
\label{sec:orgb221a3b}

The Einstein-Hilbert Lagrangian in Planck units (\(G = 1\)):

\begin{equation}
\mathcal{L}_\text{GR} = \frac{1}{16\pi} R - \Lambda
\end{equation}

where \(R\) is the Ricci scalar curvature in units of \(1/l_\text{Planck}^2\) and \(\Lambda\) is the cosmological constant in the same units. Both are dimensionless ratios in Planck units.

\subsection{The Cosmological Constant in Planck Units}
\label{sec:org08726e6}

The vacuum negotiation energy density in Planck units:

\begin{equation}
\Lambda = \frac{E_\text{Hubble}}{E_\text{Planck}} \cdot \left(\frac{l_\text{Planck}}{R_H}\right)^3 \approx 10^{-123}
\end{equation}

This is consistent in order of magnitude with the observed cosmological constant. The \(10^{123}\) hierarchy between the Planck-scale QFT prediction and the observed \(\Lambda\) is reinterpreted as the ratio of the two scale inputs --- the depth of the Medium between its grain and its horizon. Whether this reinterpretation dissolves the hierarchy problem or merely relocates it depends on whether the Hamiltonian floor at \(\rho_\Lambda\) can be independently derived (see Document 1, Section on threshold emergence). That derivation does not yet exist.

\begin{openbox}
\textbf{Open calculation 3 --- precise derivation of $\hbar$}\\[4pt]
First order: $\hbar = E_\text{Planck} \cdot t_\text{Planck} / 2\pi$ gives a value within one order of magnitude of the measured $\hbar$. The Barbero-Immirzi correction $\gamma \approx 0.2375$ modifies the minimum area quantum from $l_\text{Planck}^2$ to $\gamma \cdot l_\text{Planck}^2$, shifting $\hbar$ toward its measured value. The precise non-circular derivation yielding 99.1\% agreement is given in the companion $\hbar$ draft (see Document 1 and companion hbar section).
\end{openbox}

\section{Summary: The Complete Lagrangian}
\label{sec:orgdd611e5}

\begin{defbox}
\textbf{The complete Lagrangian in Planck units:}
\[
\mathcal{L}_\text{full} = \underbrace{\frac{1}{2}\!\left(\frac{\partial N}{\partial t}\right)^{\!2} - \frac{1}{2}\,N^\top L N}_{\text{negotiation field}} - \underbrace{\varepsilon_\text{edge}\sum_{ij}|\dot{h}_{ij}|}_{\text{topology change}} + \underbrace{\frac{1}{16\pi}R - \Lambda}_{\text{gravity}}
\]
All terms dimensionless in Planck units. Two inputs: iota size $1$ and Hubble scale $R_H$.\\[4pt]
Derived from this Lagrangian: $\hbar$, SU(3), color charge, spin, confinement, mass, $\Lambda$, gravitational coupling.
\end{defbox}

What the Lagrangian establishes:
\begin{enumerate}
\item A well-defined field theory on a graph --- the negotiation Lagrangian is a standard harmonic lattice / discrete Klein-Gordon Lagrangian. Its mathematical properties are known.
\item Bosonic normal modes --- canonical quantization of \(\mathcal{L}\) with \([N, \Pi] = i\hbar\) gives bosonic creation/annihilation operators satisfying \([\hat{a}, \hat{a}^\dagger] = 1\). The resulting statistics are Bose-Einstein.
\item Mass ratios --- the Laplacian eigenvalue spectrum determines mass ratios between modes on different graph topologies, with \(\varepsilon\) as a free scale parameter.
\item A symmetry candidate --- complex-valued fields on the triangle admit a continuous U(3) or subgroup; whether it is precisely SU(3) and whether it corresponds to QCD color is an open derivation.
\item A vacuum energy estimate --- conditional on the Hamiltonian floor condition (Document 1).
\end{enumerate}

Open problem flagged for theorists: the proton is a fermion obeying Fermi-Dirac statistics and the Pauli exclusion principle. The bosonic Lagrangian above cannot produce this. Either: (a) the negotiation field must be Grassmann-valued (anticommuting) to give fermionic statistics, or (b) a composite antisymmetrization argument must show that three bosonic quanta in the triangle mode acquire fermionic exchange statistics collectively. Neither has been established. This is a primary open problem for this description.

\section{Document Series and Cross-References}
\label{sec:org78e6a18}

This is Document 3 of 5 in The Medium series.

\begin{center}
\begin{tabular}{clp{9cm}}
\toprule
\textbf{Doc} & \textbf{File} & \textbf{Content}\\[0pt]
\midrule
0 & medium\_doc0\_introduction.org & Personal introduction; the question behind the description\\[0pt]
1 & medium\_doc1\_definition.org & The Medium: Foundation Definition (primitives, three levels)\\[0pt]
2 & medium\_doc2\_unified.org & Negotiation Modes, Particle Structure, Proton Mass\\[0pt]
3 & medium\_doc3\_lagrangian.org & \emph{This paper}\\[0pt]
4 & medium\_doc4\_hamiltonian.org & Hamiltonian, energy spectrum, canonical quantization, ADM\\[0pt]
5 & medium\_doc5\_G\_H0.org & Deriving \(G\) and \(H_0\) from substrate primitives\\[0pt]
\bottomrule
\end{tabular}
\end{center}

All documents: Zenodo DOI: 10.5281/zenodo.18927154 \(\cdot\) ORCID: 0009-0001-0328-2299 \(\cdot\) github.com/bkdepuy/The-Medium

\section{License}
\label{sec:org2734b14}

CC BY 4.0. You are free to share and adapt this material for any purpose, provided you give appropriate credit to the original author.
\end{document}
